(1) 轴上设计焦点 $(0,0,f)$ 与半径 $r_m$ 孔径点的传播路程差为 $\sqrt{f^2+r_m^2}-f$。按连续半波带定义令其等于 $m\lambda/2$，即 $\sqrt{f^2+r_m^2}=f+m\lambda/2$。平方后得到精确关系 $r_m^2=m\lambda f+m^2\lambda^2/4$；近轴近似仅在第二项远小于第一项时成立。
(2) 当 $m\lambda\ll f$ 时 $r_m^2\simeq m\lambda f$，也就是相邻半波带的半径平方差近似 $r_{m+1}^2-r_m^2\simeq\lambda f$。特别地 $r_1^2\simeq\lambda f$，所以 $f\simeq r_1^2/\lambda$。若误写成相邻半波带光程差 $\lambda$ 会造成焦距相差一倍。
(3) 相邻半波带到焦点贡献的场相位相差约 $\pi$，若把相邻带交替遮挡，只保留同奇偶性的带，它们的相位差约为 $2\pi$，可在焦点相干增强。二元周期透过率在平方半径坐标中具有奇次傅里叶级数，因此除主焦距 $f$ 外有近似焦点 $f_3=f/3,f_5=f/5,\ldots$；偶数级理想衍射级次消失，实际效率受带数与遮光方式限制。
(4) 对于固定的环带平方半径 $r_m^2\simeq m\lambda f$，换成波长 $\lambda'$ 后必须有 $r_m^2\simeq m\lambda'f'$，故主焦距 $f'\simeq f\lambda/\lambda'$。更长波长聚焦得更近，表明波带片具有强衍射色散；它与依靠玻璃折射率分散获得光焦度的折射式薄透镜不同，不具有无色散的固定焦距。
